%! TEX program = xelatex
\documentclass[a5paper, twoside, 11pt]{article}

\usepackage{enumitem, amsmath, amssymb, tikz, fancyhdr, stackengine, graphicx, sectsty, pdfpages, tabularx}

\usepackage[top=18mm, inner=10mm, outer=9mm, bottom=17mm, headheight=14.5pt, headsep=3mm, footskip=7mm]{geometry}

    \usepackage[english, russian]{babel}    %% загружает пакет многоязыковой вёрстки
    \usepackage{fontspec}                  %% подготавливает загрузку шрифтов 
    \defaultfontfeatures[\rmfamily, \sffamily]{Ligatures={TeX}}  %% свойства шрифтов по умолчанию
    \setmainfont[Scale=.9]{PT Serif} %% задаёт основной шрифт документа
    \setsansfont[Scale=.9]{PT Sans} %% задаёт шрифт без засечек
    \setmonofont{Consolas}              %% задаёт моноширинный шрифт 

    \usetikzlibrary{arrows.meta}

    \usepackage{unicode-math}
    \setmathfont{[latinmodern-math.otf]}
    \setmathfont[range=\mathit/{latin,Latin}]{PT Serif Italic}
    %\setmathfont[range=\mathit/{greek,Greek}]{Arno Pro}    
    \setmathfont[range=up]{PT Serif}
    \setmathfont[range="2264-"2265]{PT Serif}
    \setmathfont[range="003C-"003E]{PT Serif}
	
    \allsectionsfont{\sffamily}

    \usepackage{verbatim}
    \usepackage{import}

    \setlist[enumerate,1]{nolistsep}
    \setlist[itemize,1]{nolistsep}

    % русские названия заголовков разделов 
    \newcommand{\problemsname}{Задачи}
    \newcommand{\solutionsname}{Решения}
    \newcommand{\inputname}{Формат входных данных}
    \newcommand{\outputname}{Формат выходных данных}
    \newcommand{\examplename}{Пример входного и выходного файлов}
    \newcommand{\examplesname}{Примеры входного и выходного файлов}
    \newcommand{\commentsname}{Комментарии}
    \newcommand{\evaluationname}{Описание системы оценивания}
    \newcommand{\InputFileName}{input.txt}
    \newcommand{\OutputFileName}{output.txt}
    \makeatletter
    \def\kw@Explanation{Пояснение к примеру}
    \def\kw@Explanations{Пояснения к примерам}
    
    % заголовки разделов и~столбцов таблицы примеров
    \newcommand{\problempartheading}[1]{\par\medskip\noindent\textbf{\sffamily #1}\par\penalty10000\noindent}    
    \newcommand{\InputFile}{\problempartheading{\inputname}}
    \newcommand{\OutputFile}{\problempartheading{\outputname}}
    \newcommand{\Comments}{\problempartheading{\commentsname}}
    \newcommand{\Scoring}{\problempartheading{\evaluationname}}
    \newcommand{\Example}{\par\medskip\noindent\textbf{\sffamily\examplename}\raisebox{-2pt}{\strut}\par\penalty10000\noindent}
    \newcommand{\Examples}{\par\medskip\noindent\textbf{\sffamily\examplesname}\raisebox{-2pt}{\strut}\par\penalty10000\noindent}
    \newcommand{\Explanation}{\problempartheading{\kw@Explanation}}
    \newcommand{\Explanations}{\problempartheading{\kw@Explanations}}


	% русские знаки нестрогих неравенств
    \renewcommand{\geq}{\geqslant}
    \renewcommand{\leq}{\leqslant}
    \renewcommand{\ge}{\geq}
    \renewcommand{\le}{\leq}

    % штрафы за разрыв строки по оператору
    \binoppenalty=10000
    \relpenalty=10000
    
    \makeatletter
    % задача = \begin{problem}{name}{in}{out}{TL}{ML}
    \newenvironment{problem}[5]{%
\subsection{#1}%
\renewcommand{\InputFileName}{#2}%
\renewcommand{\OutputFileName}{#3}%
}{}
    \newenvironment{tutorial}[1]{\subsection{#1}}{}
    \renewcommand{\thesubsection}{\Alph{subsection}.}
    \renewcommand{\thesection}{}

    % ширина таблицы примеров - по умолчанию равна ширине текста минус 2 красных строки
    \newlength{\ex@mpwidth}
    \setlength{\ex@mpwidth}{\textwidth}
    \advance\ex@mpwidth-2\parindent
    \advance\ex@mpwidth-4\tabcolsep

    % ширины колонок ввода, вывода и~примеров для таблиц примеров
    \newlength{\ex@mpinpwidth}
    \newlength{\ex@mpoutwidth}
    \newlength{\ex@mpcomwidth}

% -- Setup sizes --
\newlength{\thelinewidth}
\thelinewidth=\textwidth
\newlength{\exmpwidinf}
\newlength{\exmpwidouf}
\newlength{\exmpwidewid}
\newlength{\exmpthreewidinf}
\newlength{\exmpthreewidouf}
\newlength{\exmpthreewidnote}

\newif\ifintentionallyblankpages

\exmpwidinf=0.43\thelinewidth
\exmpwidouf=0.43\thelinewidth
\exmpwidewid=0.9\thelinewidth
\exmpthreewidinf=0.28\thelinewidth
\exmpthreewidouf=0.28\thelinewidth
\exmpthreewidnote=0.30\thelinewidth


    % во всех примерах нужно следить за разрывами строк, вставлять знак комментария после завершения команды \exmp во избежание ненужных разрывов
    % например \exmp{2 4}{3 5}%

    
    % таблица примеров с 2 колонками (ввод и~вывод)
% :FIXME:

% This is magic, which delete space after verbatiminput
\addto@hook{\every@verbatim}{\topsep=0pt\relax}

\def\s@tm@cr@s{
    \def\widthin##1{\exmpwidinf=##1\relax}
    \def\widthout##1{\exmpwidouf=##1\relax}
    \def\stretchin##1{\advance\exmpwidinf by ##1\relax}
    \def\stretchout##1{\advance\exmpwidouf by ##1\relax}
    \@ifstar{
        \error Star must not be used in example environment any more
    }
}


\newenvironment{example}[1][]{
    \s@tm@cr@s#1
    \ttfamily\obeylines\obeyspaces\frenchspacing
    \newcommand{\exmp}[2]{
        \begin{minipage}[t]{\exmpwidinf}\rightskip=0pt plus 1fill\relax##1\medskip\end{minipage}&
        \begin{minipage}[t]{\exmpwidouf}\rightskip=0pt plus 1fill\relax##2\medskip\end{minipage}\\
        \hline
    }

    \newcommand{\exmpfile}[2]{
       \exmp{
          \verbatiminput{##1}
       }{
          \verbatiminput{##2}
       }%
    }


    \begin{tabular}{|l|l|}
        \hline
        \multicolumn{1}{|c|}{\bf\texttt{\InputFileName}}&
        \multicolumn{1}{|c|}{\bf\texttt{\OutputFileName}}\\
        \hline
}{
    \end{tabular}
}

    
    % таблица примеров с 3 колонками (ввод, вывод, комментарий), три аргумента - доли колонок 
        \newenvironment{examplecommented}[3]{
        \setlength{\ex@mpinpwidth}{#1\ex@mpwidth}
        \setlength{\ex@mpoutwidth}{#2\ex@mpwidth}
        \setlength{\ex@mpcomwidth}{#3\ex@mpwidth}
        \newcommand{\exmp}[3]{\begin{minipage}[t]{\ex@mpinpwidth}\raggedright\ttfamily ##1\medskip\end{minipage} &
                              \begin{minipage}[t]{\ex@mpoutwidth}\raggedright\ttfamily ##2\medskip\end{minipage} &
                              \begin{minipage}[t]{\ex@mpcomwidth} ##3\medskip\end{minipage}\\\hline%
        }
        \obeylines\obeyspaces\frenchspacing

        \noindent\begin{tabular}{|l|l|l|}
        \hline\multicolumn{1}{|c|}{\bf\texttt{\InputFileName}}&\multicolumn{1}{c|}{\bf\texttt{\OutputFileName}}&
        \multicolumn{1}{c|}{\bf\texttt{\commentsname}}\\
        \hline
        }{\end{tabular}\smallskip}

    
    % таблица примеров в~1 колонку. Один пример верстается в~4 строки (считая заголовки) на всю ширину таблицы
        \newenvironment{examplewide}{
        \advance\ex@mpwidth2\tabcolsep 
        \ttfamily\obeylines\obeyspaces\frenchspacing
        \newcommand{\exmp}[2]{
            \begin{tabular}{|c|}
            \hline
            \multicolumn{1}{|c|}{\bf\texttt{\InputFileName}}\\
            \hline
            \begin{minipage}[t]{\ex@mpwidth}\rightskip=0pt plus 1fill\relax\raggedright ##1 \medskip\end{minipage}\\
            \hline
            \multicolumn{1}{|c|}{\bf\texttt{\OutputFileName}}\\
            \hline
            \begin{minipage}[t]{\ex@mpwidth}\rightskip=0pt plus 1fill\relax\raggedright ##2 \medskip\end{minipage}\\%
            \hline
            \end{tabular}\smallskip
        }
        }{}

   \makeatother

	% колонтитулы
    \pagestyle{fancy}
    \makeatletter
    \fancyhead[LE]{\small \@title }
    \makeatother
    \fancyhead[RE]{}
    \fancyhead[LO]{2016--2017 учебный год}
    \fancyhead[RO]{}
    
    \newcommand\olympiad{*** ОПРЕДЕЛИТЕ ОЛИМПИАДУ КОМАНДОЙ \textbackslash olympiad !!! ***}

    
    \AtBeginDocument{\sloppy}
    \newcommand*{\hm}[1]{#1\nobreak\discretionary{}%
{\hbox{$\mathsurround=0pt #1$}}{}}


\newcommand{\foreword}{
\noindent\centerline{\parbox{0.9\textwidth}{\centering\sloppy\textbf{\olympiad}}}

\medskip
\noindent\centerline{\textbf{Дорогие участники!}}
\medskip

Проверка ваших решений производится \emph{автоматически} на наборе тестов, учитывающем различные варианты исходных данных, чтобы система могла выставить наиболее объективную оценку.

Для каждого теста выполняется следующая процедура:
\smallskip
\begin{itemize}[nolistsep]
\item
Система создает входной файл с исходными данными теста.
\item 
Запускается ваша программа, которая должна прочитать этот файл и~записать ответ в выходной файл.
\item
Затем система проверяет созданный вашей программой файл и~выносит решение о его правильности. 
\end{itemize}

Именно поэтому ваша программа должна читать данные из файла, а не с клавиатуры, и вывод записывать тоже в файл. Ни в коем случае нельзя использовать ввод с клавиатуры (например, функцию \texttt{readkey} в Паскале), так как в этом случае ваша программа будет ждать ввода бесконечно (и будет снята с тестирования после превышения лимита времени). Также важно строго соблюдать формат выходного файла. 

Приведем простой способ чтения из файла и записи в файл на языке Паскаль:\\
\framebox{\parbox{\textwidth}{
\small\tt
\{ {\rm\textit{ в начале программы }}\/ \}\\
assign(input, 'taskname.in');  reset(input);\\
assign(output, 'taskname.out'); rewrite(output);\\
\{ {\rm\textit{ теперь обычные процедуры read, readln, write, writeln будут работать \\ с файлами, а не с клавиатурой / экраном}} \}\\
  …\\
\{ {\rm\textit{ в конце программы }}\/ \}\\
close(output);
}}

\bigskip
\noindent
Ко всем задачам предъявляются следующие технические требования:

\noindent\addstackgap[8pt]{
\begin{tabular}{@{}ll}
Имена входного и выходного файлов & \textbf{см. в примерах к~задаче}\\
Максимальное время работы на одном тесте & 1 секунда\\
Максимальный объем используемой памяти	& 256 мегабайт \\
\end{tabular}
}


\begin{center}
\textbf{Желаем удачи!}
\end{center}
}


\newcommand{\importpolygontex}[2]{
\graphicspath{{problems/#1/statements/russian/}}
\import{problems/#1/statements/russian/}{#2.tex}   
}

\newcommand{\importproblem}[1]{\importpolygontex{#1}{problem}}
\newcommand{\importtutorial}[1]{\importpolygontex{#1}{tutorial}}
\usepackage{lscape,longtable,amsmath,tabularx}

\renewcommand{\olympiad}{XIX Республиканская командная олимпиада школьников \mbox{по программированию}}
\fancyhead[LE]{XIX Республиканская командная олимпиада школьников}
\fancyhead[RE]{}
\fancyhead[RO]{2021--2022 учебный год}
\fancyhead[LO]{}

\begin{document}


\pagestyle{empty}

\begin{flushright}
\it
Эффективная и слаженная\\
командная работа --- залог успеха!
\end{flushright}
{}\vskip-15mm

\includegraphics[width=3cm, height=3cm]{logo.pdf}
\\[3cm]

\begin{center}
\huge
XIX РЕСПУБЛИКАНСКАЯ \\КОМАНДНАЯ\\ОЛИМПИАДА ШКОЛЬНИКОВ \\
ПО~ПРОГРАММИРОВАНИЮ
\end{center}

\vfill
\centerline{Якутск}
\centerline{4--5 марта 2022 г.}
\newpage
\noindent XIX Республиканская командная олимпиада школьников по программированию.~--- Якутск, 2022.
\\[5mm]
Сборник содержит условия задач XIX Республиканской командной олимпиады школьников по программированию и возможные варианты решений. Олимпиада проводилась дистанционно 
4--5 марта 2022~г.
\vfill

{}\hfill © Авторский коллектив, 2022

\newpage
\noindent
\textbf{СПОНСОРЫ ОЛИМПИАДЫ}
\\[7mm]
\begin{tabular}{@{}p{.44\textwidth}p{.5\textwidth}}
\raisebox{-4.7ex}{\includegraphics{cts.png}}
& 
Генеральный директор \newline
ГК <<КопирТехСервис>>\newline
НИКИФОРОВ Александр Тарасович
\newline\strut
\\
\raisebox{-4.7ex}{\includegraphics[width=.43\textwidth]{nefu.pdf}}
& 
Преподаватели СВФУ \newline 
имени М.К.Аммосова \newline
НИКОЛАЕВА Наталья Васильевна, \newline
ЛЕВЕРЬЕВ Владимир Семенович
\newline\strut
\\
\raisebox{-4.7ex}{\includegraphics{lena.png}}
& 
Руководитель МАУ <<Кинотеатр\newline
Лена>> ГО <<город Якутск>>\newline
ЕГОРОВ Алексей Алексеевич
\end{tabular}

\newpage
\noindent
\section*{ОРГКОМИТЕТ ОЛИМПИАДЫ}
\addcontentsline{toc}{section}{Оргкомитет}
%\fancyhead[RO]{Оргкомитет олимпиады}
\thispagestyle{plain}

ГОЛИКОВ Алексей Иннокентьевич \\ 
\textit{заместитель ректора Северо-Восточного федерального \\
университета имени М.\,К. Аммосова по кадровой политике~--- \\
председатель}
\\[3mm]
ПАВЛОВ Василий Климович \\
\textit{ректор Малой академии наук РС\,(Я)~--- заместитель председателя}
\\[3mm]
ПЕТРОВА Дария Александровна \\
\textit{ведущий специалист департамента по гос. политике \\
общего образования, воспитания и дополнительного образования \\
Министерства образования и науки РС(Я)}
\\[3mm]
НИКОЛАЕВА Наталья Васильевна \\
\textit{доцент ИМИ СВФУ, заведующая кафедрой информатики \\
Малой академии наук РС\,(Я  )}
\\[3mm]
СЕМЕНОВА Галина Александровна \\
\textit{проректор по учебно-организационной и воспитательной \\ 
работе Малой академии наук РС\,(Я)} 
\\[3mm]
ОБУДОВ Владислав Антонович \\
\textit{диспетчер  Малой академии наук РС\,(Я)}
\\[3mm]

\newpage
\section*{ЖЮРИ ОЛИМПИАДЫ}
\addcontentsline{toc}{section}{Жюри}
%\fancyhead[RO]{Жюри олимпиады}
\thispagestyle{plain}

ПАВЛОВ Никифор Никитич \\
\textit{к. ф.-м. н., доцент кафедры <<Информационные технологии>>\\ ИМИ СВФУ --- председатель}
\\[2mm]
БУЛАТОВ Василий Алквиадович\\
\textit{разработчик ПО, Yandex Europe B.V.}
\\[2mm]
ГАВРИЛЬЕВ Валерий Валерьевич\\
\textit{разработчик ПО, ООО «Яндекс.Технологии»}
\\[2mm]
ЗЫКОВ Тимур Алексеевич\\ 
\textit{студент, Физтех-школа прикладной математики и информатики \\
Московского физико-технического института \\
(национального исследовательского университета)}
\\[2mm]
ИВАНОВ Виталий Витальевич\\
\textit{старший инженер-программист, ООО «Ай-Новус»}
\\[2mm]
ЛЕВЕРЬЕВ Владимир Семенович \\
\textit{старший преподаватель кафедры ИТ ИМИ СВФУ}
\\[2mm]
МАКАРОВ Денис Иванович\\
\textit{студент, факультет информационных технологий и программирования Университета ИТМО}
\\[2mm]
МЕКУМЯНОВ Семен Леонидович\\
\textit{аналитик, Pinely}
\\[2mm]
НИКИФОРОВ Дьулустан Васильевич \\
\textit{старший преподаватель кафедры ИТ ИМИ СВФУ}
\\[2mm]
НИКОЛАЕВА Наталья Васильевна \\
\textit{к. ф.-м. н., зав. кафедрой ИТ ИМИ СВФУ,\\
зав. кафедрой информатики МАН РС(Я)}
\\[2mm]
ОБУДОВ Владислав Антонович, \\
\textit{диспетчер  Малой академии наук РС\,(Я) --- секретарь жюри}
\\[2mm]
ПАВЛОВ Александр Викторович \\
\textit{к. ф.-м. н., разработчик, ИП Павлов А.В.}
\newpage\noindent
ПАРНИКОВ Василий Васильевич\\
\textit{студент, факультет информационных технологий и программирования Университета ИТМО}
\\[2mm]
СЕМЕНОВ Никита Сергеевич\\
\textit{студент, факультет компьютерных наук \\
Национального исследовательского университета \\
«Высшая школа экономики»}
\\[2mm]
ЭВЕРСТОВ Владимир Васильевич \\
\textit{старший преподаватель кафедры ИТ ИМИ СВФУ}



\newpage
% \newpage

\section*{УЧАСТНИКИ ОЛИМПИАДЫ}
\addcontentsline{toc}{section}{СПИСОК УЧАСТНИКОВ}


\begin{description}[style=unboxed, labelwidth=\linewidth, font =\bfseries, listparindent =0pt, leftmargin=15pt, parsep=0pt]
\raggedright

\item[Чурапча]

Куличкин Тимур Фёдорович,	\textit{11 кл. Чурапчинская гимназия им.~С.\,К.~Макарова,	руководитель Захаров Прокопий Прокопьевич}

Сивцев Эркин Афанасьевич,	\textit{9 кл. ЧСОШ им.~С.\,А.~Новгородова,	руководитель Прокопьев Евгений Вячеславович}

Александров Николай Васильевич,	\textit{9 кл. ЧСОШ им.~С.\,А.~Новгородова,	руководитель Прокопьев Е. В.}


\item[Нюрба НТЛ]

\textit{Нюрбинский технический лицей им.~А.\,Н.~Чусовского, \\руководитель Петров Николай Николаевич}

Федоров Альберт Викторович, \textit{9 кл.}

Александров Константин Аркадьевич, \textit{9 кл.}

Данилова Сардаана Антоновна, \textit{11 кл.}

\item[StasLug20061103 ]

\textit{10 кл., Республиканский лицей-интернат (РЛИ), \\руководитель Уваровская Мария Ивановна}

Макаров Мирон Егорович

Шамаев Баатылы Эллэй-уола

Семенов Артем Герасимович


\item[Продам гараж \quad\mdseries\textit{(11 кл. РЛИ, руководитель Уваровская М. И.)}]

Павлова Виктория Владимировна

Ефимов Степан Сергеевич

Миндаев Дамир Александрович


\item[Dream team \quad\mdseries\textit{(11 кл. РЛИ, руководитель Уваровская М. И.)}]

Дьячковский Дамир Николаевич

Петров Тимур Хрисанович

Христофоров Виталий Николаевич


\item[Эрэл-7 \quad\mdseries\textit{(11 кл. РЛИ, руководитель Титов Александр Васильевич)}]

Птицына Варвара Петровна

Степанова Яна Сергеевна

Татаринов Тимур Иванович


\item[Эрэл-89 \quad\mdseries\textit{(РЛИ, руководитель Куличкин Никифор Николаевич)}]

Павлова Аина Егоровна,	\textit{9 кл.}

Юмшанов Эрчим Петрович, \textit{8 кл.}

Решетников Ньургун Петрович, \textit{9 кл.}


\item[Сборная России \quad\mdseries\textit{(10 кл. РЛИ, руководитель Уваровская М. И.)}]

Васильев Ян Евгеньевич

Николаев Давид Андреевич

Щукин Вадим Михайлович


\item[Курочка Ряба \quad\mdseries\textit{(11 кл. РЛИ, руководитель Уваровская М. И.)}]

Фролова Ксения Георгиевна

Акимов Мир Мирович

Семенов Ефим Николаевич


\item[Авкуб \quad\mdseries\textit{(10 кл. РЛИ, руководитель Уваровская М. И.)}]

Данилова Айаана Васильевна

Ноговицын Айтал Евгеньевич

Павлова Алина Михайловна


\item[Эрэл 11-10 \quad\mdseries\textit{(РЛИ, руководитель Уваровская М. И.)}]

Николаев Владимир Аянович, 10 кл.

Крылыков Айаал Михайлович, 10 кл.

Иванова Аайа Гаврильевна, 11 кл.


\item[Команда №1 \quad\mdseries\textit{(руководитель Соловьева Ирина Васильевна)}]

Николаева Сахаайа Владимировна,	\textit{10 кл., Покровская УМГ}

Васильева Алина Алексеевна, \textit{8 кл., Покровская СОШ №3}

Петров Федот Павлович,	\textit{10 кл., Покровская СОШ №2}


\item[Команда №2]

\textit{11 кл, Октемский НОЦ, руководитель Ковров Феликс Филиппович}

Иванов Айыллаан Семенович

Васильев Руслан Николаевич

Герасимов Иннокентий Иннокентьевич


\item[Команда №3]

\textit{руководитель Мордовской Денис Андреевич}

Федоров Алексей Васильевич, \textit{7 кл., Покровская СОШ №2}

Николаева Александра Афанасьевна, \textit{7 кл., Покровская УМГ}

Никитин Егор Юрьевич, \textit{8 кл., МБНОУ Октемский НОЦ	}


\item[ЯГЛ1]

\textit{10 кл, Якутский городской лицей (ЯГЛ), \\руководитель Жиркова Марта Маратовна}

Григорьев Михаил Николаевич

Кириллин Эрхан Гаврилович

Стручков Радислав Александрович


\item[ЯГЛ2 \quad\mdseries\textit{(9 кл, ЯГЛ, руководитель Жиркова М. М.)}]

Ксенофонтов Артем Александрович

Колодезников Андрей Аркадьевич

Амвросьев Данил Николаевич


\item[ФТЛ11]

\textit{11 кл., Физико-технический лицей имени В.П.Ларионова (ФТЛ),\\руководитель Куличкин Н. Н.}

Захаров Айтал Анатольевич

Белоусов Иван Алексеевич

Павловский Владислав Витальевич


\item[ФТЛ9 \quad\mdseries\textit{(9 кл., ФТЛ, руководитель Куличкин Н. Н.)}]

Романюк Анастасия Руслановна

Дмитриева Александра Степановна

Кривошапкин Эдуард Алексеевич


\item[ФТЛ10 \quad\mdseries\textit{(10 кл., ФТЛ, руководитель Куличкин Н. Н.)}]

Местников Марат Павлович

Привалов Владимир Денисович

Павлов Даниель Альбертович


\item[СПЛ-1 \quad\mdseries\textit{(Саха политехнический лицей (СПЛ))}]

Кузнецов Юрий Данилович, \textit{11 кл., руководитель Герасимова Галина Егоровна}

Фёдоров Артём Афанасьевич, \textit{9 кл., руководитель Саввинов Игорь Семенович}

Егорова Лилия Михайловна, \textit{10 кл., руководитель Фролова Светлана Марковна}


\item[СПЛ-2 \quad\mdseries\textit{(10 кл. СПЛ)}]

Кондратьев Анатолий Львович, \textit{руководитель Герасимова Г.\,Е.}

Павлов Денис Иванович, \textit{руководитель Саввинов И.\,С.}

Борисов Никита Александрович, \textit{руководитель Фролова С.\,М.}


\item[СПЛ-3 \quad\mdseries\textit{(7 кл. СПЛ)}]

Нектегяев Марк Васильевич, \textit{руководитель Герасимова Г.\,Е.}

Петров Кирилл Владимирович, \textit{руководитель Саввинов И.\,С.}

Яковлев Арсен Викторович, \textit{руководитель Фролова С.\,М.}


\item[RumpleThump]

\textit{10 кл., средняя школа № 26 им. Е.\,Ю. Келле-Пелле г. Якутска, \\руководитель Алексеев Алексей Александрович}

Дьячковский Петр Владиславович

Погодаев Александр Михайлович

Прокопьева Глория Егоровна


\item[СОШ9]

\textit{средняя школа №9 им. М.\,И. Кершенгольца г. Якутска, \\руководитель Слепцов Николай Васильевич}

Исаев Андрей Васильевич, \textit{9 кл.}

Шишкин Кирилл Денисович, \textit{9 кл.}

Аркадьева Ксения Николаевна, \textit{10 кл.}


\item[10-МИТ]

\textit{10 кл., Специализированный учебно-научный центр --- \\Университетский лицей СВФУ (СУНЦ), \\руководитель Васильева Наталья Васильевна}

Слепцов Никита Сергеевич

Попов Кирилл Константинович

Федоров Богдан Иванович


\item[11-МИТ \quad\mdseries\textit{(11 кл., СУНЦ СВФУ, руководитель Васильева Н. В.)}]

Ефимов Владимир Михайлович

Тарасенко Антон Васильевич

Кардашевский Ян Михайлович


\item[МАШ]

\textit{Международная Арктическая школа (МАШ), \\руководитель Ильин Максим Макарович}

Дуткин Тимур Евгеньевич, \textit{11 кл.}

Мыреева Ирина Алексеевна, \textit{11 кл.}

Сергин Сергей Викторович, \textit{10 кл.}


\item[Analogy \quad\mdseries\textit{(9 кл. МАШ, руководитель Ильин М. М.)}]

Петров Леонид Николаевич

Гаврильева Камилла Максимовна

Капитонов Денис Евгеньевич
\end{description}

\pagestyle{fancy}

\newpage
\rm
\thispagestyle{plain}
\section*{УСЛОВИЯ ЗАДАЧ}
\addcontentsline{toc}{section}{Условия задач}
\fancyhead[RO]{Условия задач}

\newcommand{\importproblem}[1]{
\graphicspath{{problems/#1/statements/russian/}}
\import{problems/#1/statements/russian/}{problem.tex}	
}

\importproblem{thereandback}
\importproblem{xor-from-a-to-b}
\importproblem{pizza}
\importproblem{dyyray-army}
\importproblem{dna}
\importproblem{pelmeni}
\importproblem{kihun}
\importproblem{circles}
\importproblem{containers}
\importproblem{fractions}


%======================================================================
\newpage
\fancyhead[RO]{Решения}
\section*{РЕШЕНИЯ}
\addcontentsline{toc}{section}{Решения}
\setcounter{subsection}{0}


\renewcommand{\importproblem}[1]{
\graphicspath{{problems/#1/statements/russian/}}
\import{problems/#1/statements/russian/}{tutorial.tex}	
}

\importproblem{thereandback}
\importproblem{xor-from-a-to-b}
\importproblem{pizza}
\importproblem{dyyray-army}
\importproblem{dna}
\importproblem{pelmeni}
\importproblem{kihun}
\importproblem{circles}
\newpage
\importproblem{containers}
\importproblem{fractions}


\begin{landscape}
% \documentclass[a5paper, 10pt, landscape]{book}
% \usepackage[hmargin=17mm,vmargin=10mm]{geometry}
% \usepackage{fontspec,longtable,tabularx,amsmath}
% \setmainfont{PT Sans}
% \setsansfont{PT Sans}

% \begin{document}
% \pagestyle{empty}
%
% \newcommand{\ptim}[2]{$\underset{\mbox{\sffamily\scriptsize #2}}{\mbox{#1}}$}
\sffamily
\pagestyle{empty}
\setlength{\tabcolsep}{.7mm}
\sectionfont{\nohang\centering\sffamily}
\section*{ПРОТОКОЛ}
\addcontentsline{toc}{section}{Итоги олимпиады}
\vskip-3mm
\centerline{XX Республиканская командная олимпиада школьников по программированию}
\fontsize{10pt}{12pt}\selectfont\par\noindent
\begin{longtable}{
|r|
>{\raggedright\arraybackslash}m{40mm}|
*{11}{>{\centering\arraybackslash}m{7mm}|}
>{\centering\arraybackslash\bfseries\sffamily}m{7mm}|
>{\raggedleft\small}m{9.5mm}|
>{\centering\arraybackslash\small}m{12mm}|
}
\hline
№ & Команда & A & B & C & D & E & F & G & H & I & J & K & = & {\small Время} & Диплом
\tabularnewline
\hline
1 & \textbf{Чурапча} \linebreak (Куличкин Т.Ф., Сивцев Э.А., Александров Н.В.) & +1 & +2 & + & + & + & +2 & + & +2 & 1 & + & +3 & 11 & 1126 & I степени \\ \hline
2 & \textbf{ФТЛ11}, Якутск (Захаров А.А., Белоусов И.А., Павловский В.В.) & + & +1 & + & +3 & +1 & +2 & + & +1 & + & + & +3 & 11 & 1269 & I степени \\ \hline
3 & \textbf{StasLug20061103}, РЛИ (Макаров М.Е., Шамаев Б.Э., Семенов А.Г.) & + & + & + & +7 & +2 & +3 & + & +2 & + & + & +3 & 11 & 1503 & I степени \\ \hline
4 & \textbf{Продам гараж}, РЛИ (Павлова В.В., Ефи- \linebreak мов С.С., Миндаев Д.А.) & + & +1 & +1 & +1 & +1 & +3 & + & +1 & + & + & + & 11 & 1552 & I степени \\ \hline
5 & \textbf{ФТЛ9}, Якутск (Романюк А.Р., Дмитриева А.С., Кривошапкин Э.А.)  & + & +3 & + & +1 & +1 & +1 & + & -3 &  & + &  & 8 & 984 & II степени \\ \hline
6 & \textbf{Курочка} Ряба, РЛИ (Фролова К.Г., Аки-\linebreak мов М.М., Семенов Е.Н.) & + & +1 & + & + &  & + & + & -3 &  & +1 & -1 & 7 & 461 & II степени \\ \hline
7 & \textbf{Dream} team, РЛИ (Дьяч-\linebreak ковский Д.Н., Петров Т.Х., Христофоров В.Н.) & + & + & + & +4 &  & +1 & + & -2 &  & + &  & 7 & 738 & II степени \\ \hline
8 & \textbf{ЯГЛ2}, Якутск (Ксенофонтов А.А., Колодезников А.А., Амвросьев Д.Н.) & + & + & + &  &  & +6 & + & +2 & -2 & + &  & 7 & 754 & II степени \\ \hline
9 & \textbf{Сборная} России, РЛИ (Васильев Я.Е., Николаев Д.А., Щукин В.М.) & +1 & +2 & +1 & +3 &  & +10 & +1 & -7 &  & + &  & 7 & 913 & II степени \\ \hline
10 & \textbf{СПЛ-1}, Якутск (Кузнецов Ю.Д., Фёдоров А.А., Егорова Л.М.) & + & + & + &  &  & +2 & + &  &  & + &  & 6 & 269 & III степени \\ \hline
11 & \textbf{ФТЛ10}, Якутск (Местников М.П., Привалов В.Д., Павлов Д.А.) & + & +1 & + & +4 &  & -5 & + & -8 &  & + &  & 6 & 337 & III степени \\ \hline
12 & \textbf{Эрэл-89}, РЛИ (Павлова А.Е., Юмшанов Э.П., Решетников Н.П.) & + & +2 & + & +2 &  & -5 & + &  &  & + &  & 6 & 479 & III степени \\ \hline
13 & \textbf{Авкуб}, РЛИ (Данилова А.В., Ноговицын А.Е., Павлова А.М.) & + & +1 & +1 & +2 &  & -1 & + & -1 &  & + &  & 6 & 575 & III степени \\ \hline
14 & \textbf{ЯГЛ1}, Якутск (Григорьев М.Н., Кириллин Э.Г., Стручков Р.А.) & + & + & + &  &  & -1 & + & -6 &  & + &  & 5 & 101 & Похв.гр. \\ \hline
15 & \textbf{Эрэл 11-10}, РЛИ (Николаев В.А., Крылыков А.М., Иванова А.Г.) & +1 & +1 & + & -3 &  & -3 & + & -2 & -5 & +1 &  & 5 & 199 & Похв.гр. \\ \hline
16 & \textbf{МАШ} (Дуткин Т.Е., Мыреева И.А., Сергин С.В.) & + & + & + &  &  & -2 & + &  &  & + &  & 5 & 312 & Похв.гр. \\ \hline
17 & \textbf{СОШ9}, Якутск (Исаев А.В., Шишкин К.Д., Аркадьева К.Н.) & + & +2 & + & -2 &  &  & + & -2 &  & + &  & 5 & 344 & Похв.гр. \\ \hline
18 & \textbf{RumpleThump}, СОШ26, Якутск (Дьячковский П.В., Погодаев А.М., Прокопьева Г.Е.) & + & +2 & +2 &  &  &  & + &  &  & + & -2 & 5 & 401 & Похв.гр. \\ \hline
19 & \textbf{11-МИТ}, СУНЦ СВФУ (Ефимов В.М., Тарасенко А.В., Кардашевский Я.М.) & + & +5 & +2 & -1 & -1 &  & + & -3 & -1 & +3 & -1 & 5 & 649 & Похв.гр. \\ \hline
20 & \textbf{СПЛ-2}, Якутск (Кондратьев А.Л., Павлов Д.И., Борисов Н.А.) & + & -4 & + &  &  & -4 & +1 &  &  & +3 &  & 4 & 313 &  \\ \hline
21 & \textbf{Команда №2}, Покровск/Октемцы (Иванов А.С., Васильев Р.Н., Герасимов И.И.) & + & -2 & + &  &  & -1 & +2 & -2 &  & + & -3 & 4 & 366 &  \\ \hline
22 & \textbf{Нюрба НТЛ} (Федоров А.В., Александров К.А., Данилова С.А.) & + & -20 & + &  &  &  &  & -6 &  & + &  & 3 & 113 &  \\ \hline
23 & \textbf{Команда №1}, Покровск/Октемцы (Николаева С.В., Васильева А.А., Петров Ф.П.) & + &  &  & -1 &  &  & +2 &  &  & + &  & 3 & 340 &  \\ \hline
24 & \textbf{Эрэл-7}, РЛИ (Птицына В.П., Степанова Я.С., Татаринов Т.И.) & + &  &  &  &  &  & + &  &  & -10 &  & 2 & 73 &  \\ \hline
25 & \textbf{Analogy}, МАШ (Петров Л.Н., Гаврильева К.М., Капитонов Д.Е.) & + &  &  &  &  &  & -2 &  &  & +2 &  & 2 & 127 &  \\ \hline
26 & \textbf{СПЛ-3}, Якутск (Нектегяев М.В., Петров К.В., Яковлев А.В.) & + &  &  &  &  &  &  &  &  & +4 &  & 2 & 167 &  \\ \hline
27 & \textbf{Команда №3}, Покровск/Октемцы (Федоров А.В., Николаева А.А., Никитин Е.Ю.) & + &  &  &  &  &  &  &  &  & +8 &  & 2 & 312 &  \\ \hline
\multicolumn{16}{|l|}{вне зачета:}\\ \hline
--- & \textbf{10-МИТ}, СУНЦ СВФУ (Попов К.К., Федоров Б.И.)  & + &  & -4 &  &  & -5 &  &  &  & + &  & 2 & 60 & \\ \hline
\end{longtable}

\smallskip\noindent Председатель жюри: Н.\,Н. Павлов\hfill 4 марта 2023 г.

\newpage
%\end{document}
\end{landscape}

\tableofcontents
\end{document}